\documentclass[10pt,a4paper]{article}
\usepackage[utf8]{inputenc}
\usepackage[T1]{fontenc}
\usepackage{lmodern}
\usepackage[brazilian]{babel}
\usepackage{amsmath,amssymb}
\usepackage[margin=1.2cm,top=1.0cm,bottom=1.0cm]{geometry}
\usepackage{enumitem}
\usepackage{xcolor}
\usepackage{tcolorbox}

\pagestyle{empty}

\tcbset{
    colback=red!5!white,
    colframe=red!75!black,
    fonttitle=\bfseries\small,
    title=ADVERTÊNCIA OBRIGATÓRIA E INSTRUÇÕES GERAIS,
    boxrule=0.7pt,
    arc=1.5mm,
    left=2mm, right=2mm, top=1.5mm, bottom=1.5mm
}

\begin{document}

% =============================================================================
% PÁGINA 1 (FRENTE)
% =============================================================================

% --- CABEÇALHO COMPACTO ---
\noindent
\begin{minipage}[t]{0.70\textwidth}
    \textbf{\small UNIVERSIDADE FEDERAL DE SANTA CATARINA -- DEPTO. DE FÍSICA}\\
    \textbf{\small FSC5101 -- Física I $\vert$ Prova: Capítulos 4 e 5 (Sears \& Zemansky)}\\
    \textbf{\small Aluno(a):} \underline{\hspace{8.2cm}}
\end{minipage}
\hfill
\begin{minipage}[t]{0.28\textwidth}
    \raggedleft
    \textbf{\small Data:} \underline{\hspace{2.2cm}}\\[1mm]
    \textbf{\small Turma:} \underline{\hspace{1.8cm}}\\[1mm]
    \textbf{\small Matrícula:} \underline{\hspace{2.0cm}}
\end{minipage}

\vspace{1.5mm}

% --- CAIXA DE ADVERTÊNCIA E INSTRUÇÕES ---
\begin{tcolorbox}
\footnotesize
\textbf{1. Proibição Estrita de IA / LLMs:} É terminantemente proibido o uso de Inteligência Artificial, geradores de texto e Modelos de Linguagem (LLMs). As soluções serão auditadas de forma independente por \textbf{3 modelos de IA distintos}; caso \textbf{2 ou mais modelos} classifiquem uma ou mais resoluções como suspeitas de geração automatizada, a prova será anulada e o aluno receberá \textbf{NOTA ZERO automaticamente}.\\
\textbf{2. Exigência de DCL e Leis da Física:} Em todas as questões é \textbf{OBRIGATÓRIO} desenhar o \textbf{Diagrama de Corpo Livre (DCL)} identificando cada força atuante e declarar explicitamente \textbf{qual lei da física} (1ª, 2ª ou 3ª Lei de Newton) está sendo aplicada em cada componente ($x$ e $y$). Respostas sem DCL ou sem justificativa da lei não serão pontuadas.\\
\textbf{3. Dados:} Adote $g = 9,80\text{ m/s}^2$ em todos os cálculos. Todas as 5 questões têm peso igual ($2,0$ pontos cada). Responda a caneta nos espaços demarcados.
\end{tcolorbox}

\vspace{1.5mm}

% --- QUESTÃO 1 (FÁCIL - SEARS 4.24) ---
\noindent
\textbf{1. [Fácil -- 2,0 pontos] [Sears 4.24]} O piso de um elevador exerce uma força normal de $620\text{ N}$ de baixo para cima sobre um passageiro que pesa $650\text{ N}$.
\begin{enumerate}[label=\alph*),itemsep=0pt,topsep=1pt]
    \item Quais são as reações dessas duas forças no sentido da Terceira Lei de Newton (especifique o corpo sobre o qual cada reação atua, sua direção e sentido)?
    \item O passageiro está sendo acelerado? Em caso afirmativo, faça o DCL e determine o módulo, a direção e o sentido da sua aceleração.
\end{enumerate}

\vspace{0.15cm}
\noindent\textit{\footnotesize Espaço para DCL e Resolução da Questão 1:}
\vspace{4.2cm}
\hrule
\vspace{2mm}

% --- QUESTÃO 2 (MÉDIA - SEARS 5.23) ---
\noindent
\textbf{2. [Média -- 2,0 pontos] [Sears 5.23]} Um tobogã com massa de $75,0\text{ kg}$ desce uma encosta coberta de neve inclinada de $25,0^\circ$ com a horizontal. O coeficiente de atrito cinético entre os patins do tobogã e a neve é $\mu_c = 0,12$.
\begin{enumerate}[label=\alph*),itemsep=0pt,topsep=1pt]
    \item Desenhe o DCL do tobogã com os eixos coordenados apropriados (paralelo e perpendicular à rampa).
    \item Calcule o módulo da força normal exercida pela rampa sobre o tobogã, explicitando a lei aplicada no eixo perpendicular.
    \item Calcule o módulo da aceleração do tobogã ao longo da rampa, indicando a lei aplicada no eixo paralelo.
\end{enumerate}

\vspace{0.15cm}
\noindent\textit{\footnotesize Espaço para DCL e Resolução da Questão 2:}
\vspace{4.5cm}
\hrule
\vspace{2mm}

% --- QUESTÃO 3 (MÉDIA - SEARS 5.19) ---
\noindent
\textbf{3. [Média -- 2,0 pontos] [Sears 5.19]} \textbf{Máquina de Atwood.} Uma carga de tijolos com massa $m_1 = 15,0\text{ kg}$ está suspensa por uma corda leve e inextensível que passa por uma polia fixa sem atrito. Um contrapeso de massa $m_2 = 28,0\text{ kg}$ está preso na outra extremidade. O sistema é liberado a partir do repouso.
\begin{enumerate}[label=\alph*),itemsep=0pt,topsep=1pt]
    \item Desenhe o DCL separado para cada um dos dois corpos ($m_1$ e $m_2$), indicando os sentidos assumidos para a aceleração.
    \item Escreva a Segunda Lei de Newton para cada corpo e determine o módulo da aceleração do sistema.
    \item Calcule a força de tração $T$ na corda durante o movimento acelerado.
\end{enumerate}

\vspace{0.15cm}
\noindent\textit{\footnotesize Espaço para DCL e Resolução da Questão 3:}
\vspace{4.5cm}

\newpage
% =============================================================================
% PÁGINA 2 (VERSO)
% =============================================================================

\noindent
\textbf{UFSC -- FSC5101 Física I $\vert$ Avaliação Capítulos 4 e 5 (Continuação -- Verso)} \hfill \textbf{Aluno(a):} \underline{\hspace{5.5cm}}

\vspace{2mm}
\hrule
\vspace{3mm}

% --- QUESTÃO 4 (DIFÍCIL - SEARS 5.71) ---
\noindent
\textbf{4. [Difícil -- 2,0 pontos] [Sears 5.71]} Dois blocos retangulares de massas $m_A = 2,0\text{ kg}$ e $m_B = 4,0\text{ kg}$ estão sobrepostos: o bloco A está apoiado diretamente sobre o topo do bloco B. O coeficiente de atrito estático entre as superfícies em contato dos blocos A e B é $\mu_e = 0,30$. A superfície da mesa horizontal sobre a qual o bloco B se apoia é perfeitamente lisa (sem atrito). Uma força horizontal $F$ é aplicada diretamente sobre o bloco inferior B.
\begin{enumerate}[label=\alph*),itemsep=0pt,topsep=1pt]
    \item Desenhe o DCL individual para o bloco A e para o bloco B, identificando claramente o par de ação e reação que atua entre eles.
    \item Qual é a aceleração horizontal máxima que o bloco A pode experimentar antes de começar a escorregar sobre o bloco B? Qual força é responsável por acelerar o bloco A?
    \item Qual é o módulo máximo da força horizontal $F$ que pode ser aplicada ao bloco B para que ambos os blocos acelerem juntos sem que haja deslizamento relativo entre eles?
\end{enumerate}

\vspace{0.2cm}
\noindent\textit{\footnotesize Espaço para DCL e Resolução da Questão 4:}
\vspace{9.5cm}
\hrule
\vspace{3mm}

% --- QUESTÃO 5 (DIFÍCIL - SEARS 5.77) ---
\noindent
\textbf{5. [Difícil -- 2,0 pontos] [Sears 5.77]} Um bloco com massa $m = 3,00\text{ kg}$ está sobre uma rampa inclinada de $37,0^\circ$ com a horizontal ($\sin 37,0^\circ = 0,60$; $\cos 37,0^\circ = 0,80$). O coeficiente de atrito cinético entre o bloco e a rampa é $\mu_c = 0,30$ e o coeficiente de atrito estático é $\mu_e = 0,40$. O bloco é lançado para cima ao longo da rampa com velocidade inicial $v_0 = 8,00\text{ m/s}$.
\begin{enumerate}[label=\alph*),itemsep=0pt,topsep=1pt]
    \item Desenhe o DCL do bloco durante a subida. Escreva as equações da 1ª e 2ª Leis nos eixos perpendicular e paralelo, e calcule a distância $d$ percorrida rampa acima até o bloco parar momentaneamente.
    \item No ponto mais alto, faça a análise de forças e determine se o bloco permanece em repouso no topo ou se desce de volta. Em caso de descida, desenhe o novo DCL e determine a velocidade $v$ com que o bloco atinge a base da rampa.
\end{enumerate}

\vspace{0.2cm}
\noindent\textit{\footnotesize Espaço para DCL e Resolução da Questão 5:}
\vspace{9.5cm}

\end{document}
